\chapter{Representation theory of $\mathfrak{su}(n)$}

\section{Representation theory of $\mathfrak{su}(2)$ and $\mathfrak{sl}(2)$}

\rmkb{
    Recall that $\mathrm{SU}(2)$ and $\mathrm{SO}(3)$ share the same Lie algebra (via the double
    cover $\mathrm{SU}(2)\twoheadrightarrow\mathrm{SO}(3)$):
    \[
        \mathfrak{su}(2)
        =
        \{ A \in \mathrm{Mat}(2,\mathbb{C}) \mid \mathrm{tr}(A)=0,\ A^\dagger = -A \}.
    \]
    The Pauli matrix
    \[
        \sigma_x = \begin{pmatrix}0&1\\1&0\end{pmatrix}
    \]
    is Hermitian, so $\sigma_x^\dagger = \sigma_x \neq -\sigma_x$, and hence
    $\sigma_x\notin\mathfrak{su}(2)$. The associated basis element of $\mathfrak{su}(2)$ is
    $X := i\sigma_x$, and one has $\sigma_x = -iX$.
}

\defn[complexification]{Complexification of a Lie algebra}{
    Let $\mathfrak{g}$ be a real Lie algebra. Its \textbf{complexification}
    is the complex vector space
    \[
        \mathfrak{g}_\mathbb{C} := \mathfrak{g} \otimes_\mathbb{R} \mathbb{C},
    \]
    equipped with the complex-linear extension of the Lie bracket, defined on pure tensors by
    \[
        [X \otimes \lambda,\; Y \otimes \lambda']
        :=
        [X,Y] \otimes \lambda\lambda'
        \qquad
        \forall\, X,Y \in \mathfrak{g},\ \lambda,\lambda' \in \mathbb{C}.
    \]
}

\ex{
    The complexification of $\mathfrak{su}(2)$ is
    \[
        \mathfrak{sl}(2,\mathbb{C})
        =
        \{ A \in \mathrm{Mat}(2,\mathbb{C}) \mid \mathrm{tr}(A)=0 \}.
    \]
}

\rmkb{
    Let $\mathfrak{g}$ be a real Lie algebra. Complex representations of $\mathfrak{g}_\mathbb{C}$,
    i.e.\ complex-linear Lie algebra homomorphisms $\mathfrak{g}_\mathbb{C}\to\mathrm{End}(V)$,
    are in bijection with real Lie algebra homomorphisms $\mathfrak{g}\to\mathrm{End}(V)$.
    The bijection is given by restriction to $\mathfrak{g}\subset\mathfrak{g}_\mathbb{C}$,
    with inverse given by complex-linear extension.
}

To study the representations of $\mathfrak{su}(2)$, we use this bijection and work instead
with the complex representations of $\mathfrak{sl}(2,\mathbb{C})$. We choose the standard basis
\begin{align*}
    H &:= \begin{pmatrix}1&0\\0&-1\end{pmatrix} = \sigma_z = L_z,\\[4pt]
    E &:= \begin{pmatrix}0&1\\0&0\end{pmatrix} = \tfrac{1}{2}(\sigma_x + i\sigma_y) = L_+,\\[4pt]
    F &:= \begin{pmatrix}0&0\\1&0\end{pmatrix} = \tfrac{1}{2}(\sigma_x - i\sigma_y) = L_-.
\end{align*}
One verifies directly that the commutation relations are
\[
    [H,E] = 2E, \qquad [H,F]=-2F, \qquad [E,F]=H.
\]

\propp{
    Let $\rho:\mathfrak{sl}(2,\mathbb{C})\to\mathrm{End}(V)$ be a complex representation, and let
    $v\in V$ be an eigenvector of $\rho(H)$ with eigenvalue $\lambda\in\mathbb{C}$:
    \[
        \rho(H)v = \lambda v.
    \]
    Then $\rho(E)v$ is either zero or an eigenvector of $\rho(H)$ with eigenvalue $\lambda+2$,
    and $\rho(F)v$ is either zero or an eigenvector of $\rho(H)$ with eigenvalue $\lambda-2$.
}{
    Since $\rho$ is a Lie algebra homomorphism,
    \[
        \rho(H)\rho(E)v - \rho(E)\rho(H)v = \rho([H,E])v = \rho(2E)v = 2\rho(E)v.
    \]
    Hence
    \[
        \rho(H)\rho(E)v
        =
        \rho(E)\rho(H)v + 2\rho(E)v
        =
        \lambda\rho(E)v + 2\rho(E)v
        =
        (\lambda+2)\rho(E)v.
    \]
    The same computation using $[H,F]=-2F$ gives
    \[
        \rho(H)\rho(F)v = (\lambda-2)\rho(F)v.
    \]
}

Thus $\rho(E)$ raises and $\rho(F)$ lowers the $\rho(H)$-eigenvalue by $2$. Denote by
\[
    V_\mu := \ker\!\bigl(\rho(H) - \mu\,\mathrm{id}_V\bigr)
\]
the $\mu$-eigenspace of $\rho(H)$. Starting from a single eigenvector $v\in V_\lambda$, let $V'$
be the subrepresentation generated by acting with $\rho(E)$ and $\rho(F)$ on $v$. Then
\[
    V' = \bigoplus_{n\in\mathbb{Z}} V_{\lambda+2n}.
\]
If $V$ is finite-dimensional, only finitely many summands $V_{\lambda+2n}$ are non-zero. In
particular, there exists a maximal eigenvalue $\lambda_{\max}$ such that $\rho(E)$ annihilates
$V_{\lambda_{\max}}$:

\begin{center}
\begin{tikzpicture}[>=Stealth, baseline=(current bounding box.center)]
    \node at (-1.5, 0) {$\cdots$};
    \filldraw (0,0) circle (2pt);
    \node[below=6pt] at (0,0) {$V_{\lambda_{\max}-2}$};
    \filldraw (4,0) circle (2pt);
    \node[below=6pt] at (4,0) {$V_{\lambda_{\max}}$};
    \draw[->, bend left=25] (0.15, 0.08) to node[above] {\small$E$} (3.85, 0.08);
    \draw[->, bend left=25] (3.85,-0.08) to node[below] {\small$F$} (0.15,-0.08);
    \draw[->] (-0.9, 0) -- (-0.15, 0);
    \draw[->] (4.15, 0) -- node[above] {\small$E$} (5.2, 0);
    \node at (5.5, 0) {$0$};
\end{tikzpicture}
\end{center}

\defn[-]{Highest weight vector}{
    Let $\rho:\mathfrak{sl}(2,\mathbb{C})\to\mathrm{End}(V)$ be a complex representation.
    A non-zero vector $v\in V$ is called a \textbf{highest weight vector}\index{weight!highest weight vector}
    if
    \[
        \rho(E)v = 0.
    \]
    Its $\rho(H)$-eigenvalue is called the \textbf{highest weight}\index{weight!highest weight} of $v$.
}

\propp{
    Let $V = \bigoplus_{n \in \mathbb{Z}} V_{\lambda+2n}$ be an irreducible complex
    representation of $\mathfrak{sl}(2,\mathbb{C})$, and let $v \in V_\alpha$ with
    $\alpha = \lambda_{\max}$ be a highest weight vector. Then $V$ is generated as a
    vector space by $\{\rho(F)^k v\}_{k \in \mathbb{N}_0}$.
}{
    Let $W := \operatorname{span}_\mathbb{C}\{\rho(F)^k v \mid k \in \mathbb{N}_0\} \subseteq V$.
    We show that $W$ is a subrepresentation by checking invariance under each of
    $\rho(F)$, $\rho(H)$, and $\rho(E)$.
    \medskip
    Invariance under $\rho(F)$:
    $$\rho(F)\cdot\rho(F)^k v = \rho(F)^{k+1}v \in W$$
    \medskip
    Invariance under $\rho(H)$. Since $\rho(F)^k v \in V_{\alpha-2k}$
    $$\rho(H)\rho(F)^k v = (\alpha-2k)\rho(F)^k v \in W$$
    \medskip
    Invariance under $\rho(E)$. Since $[E,F]=H$, we have
    $$\rho(E)\rho(F) = \rho(F)\rho(E) + \rho(H)$$
    For $k=1$:
    \[
        \rho(E)\rho(F)v = \rho(F)\rho(E)v + \rho(H)v = 0 + \alpha v = \alpha v.
    \]
    For general $k$, apply this identity repeatedly to expand $\rho(E)\rho(F)^k v$:
    \begin{align*}
        \rho(E)\rho(F)^k v
        &= \rho(F)\rho(E)\rho(F)^{k-1}v + \rho(H)\rho(F)^{k-1}v \\
        &= \rho(F)^2\rho(E)\rho(F)^{k-2}v
           + \rho(F)\rho(H)\rho(F)^{k-2}v
           + (\alpha-2k+2)\rho(F)^{k-1}v \\
        &= \;\cdots \\
        &= \rho(F)^k\rho(E)v
           + \bigl(\alpha + (\alpha-2) + \cdots + (\alpha-2k+2)\bigr)\rho(F)^{k-1}v.
    \end{align*}
    Since $\rho(E)v = 0$ and the sum has $k$ terms,
    \[
        \rho(E)\rho(F)^k v
        =
        k(\alpha - k + 1)\,\rho(F)^{k-1}v
        \;\in\; W.
    \]

    Thus $W$ is a subrepresentation. Since $v \neq 0$ we have $W \neq \{0\}$, and since
    $V$ is irreducible we conclude $W = V$.
}

\rmkb{
    The proof above also establishes the formula
    \[
        \rho(E)\rho(F)^k v = k(\alpha-k+1)\,\rho(F)^{k-1}v
        \qquad
        \forall\, k \in \mathbb{N}_0.
    \]
}

\corp{
    All weight spaces $V_\kappa$ are $1$-dimensional.
}{
    The vectors $v,\rho(F)v,\ldots,\rho(F)^{\alpha}v$ form a basis of $V$ (as shown in
    the proof of the following corollary). Since $\rho(F)^k v \in V_{\alpha-2k}$ and the
    eigenvalues $\alpha, \alpha-2, \ldots, -\alpha$ are pairwise distinct, each eigenspace
    is spanned by exactly one basis vector, hence $1$-dimensional.
}{}

\corp{
    The highest weight satisfies $\alpha \in \mathbb{N}_0$, and $\dim V = \alpha+1$.
}{
    Since $V$ is finite-dimensional and the non-zero vectors among $\{\rho(F)^k v\}_{k\geq 0}$
    are linearly independent (they lie in distinct eigenspaces $V_{\alpha-2k}$ of $\rho(H)$),
    there exists a smallest $k_1 \in \mathbb{N}$ with $\rho(F)^{k_1}v = 0$.

    Since $\rho(F)^{k_1}v = 0$, we have $\rho(E)\rho(F)^{k_1}v = 0$. Applying the formula
    from the remark gives
    \[
        0
        =
        k_1(\alpha - k_1 + 1)\,\rho(F)^{k_1-1}v.
    \]
    By minimality of $k_1$, the vector $\rho(F)^{k_1-1}v \neq 0$. Since $k_1 \geq 1$,
    both $k_1$ and $\rho(F)^{k_1-1}v$ are non-zero, so we must have
    \[
        \alpha - k_1 + 1 = 0
        \qquad\Rightarrow\qquad
        \alpha = k_1 - 1 \in \mathbb{N}_0.
    \]
    The vectors $v, \rho(F)v, \ldots, \rho(F)^{k_1-1}v$ are linearly independent and span
    $V$ by the previous proposition, so $\dim V = k_1 = \alpha + 1$.
}{}

\ex{
    \begin{itemize}
        \item $\mathfrak{sl}(2,\mathbb{C})$ has a $1$-dimensional representation ($\alpha=0$):
        \[
            \rho : \mathfrak{sl}(2,\mathbb{C}) \to \mathrm{End}(\mathbb{C}),
            \qquad X \mapsto 0.
        \]
        The single weight space is $V_0 = \mathbb{C}$:
        \begin{center}
        \begin{tikzpicture}[>=Stealth, baseline=(current bounding box.center)]
            \filldraw (0,0) circle (2pt) node[below=5pt] {$V_0$};
        \end{tikzpicture}
        \end{center}

        \item $\mathfrak{sl}(2,\mathbb{C})$ has a $2$-dimensional representation ($\alpha=1$), the
        \textbf{standard representation}\index{standard representation}:
        \[
            \rho : \mathfrak{sl}(2,\mathbb{C}) \to \mathrm{End}(\mathbb{C}^2),
            \qquad X \mapsto X.
        \]
        Let $x,y$ be a basis of $\mathbb{C}^2$ with $\rho(H)x = x$ and $\rho(H)y = -y$.
        Then $V_1 = \mathbb{C}x$ and $V_{-1} = \mathbb{C}y$:
        \begin{center}
        \begin{tikzpicture}[>=Stealth, baseline=(current bounding box.center)]
            \filldraw (-2.5,0) circle (2pt);
            \node[below=5pt] at (-2.5,0) {$V_{-1} = \mathbb{C}y$};
            \node[above=5pt] at (-2.5,0) {$-1$};
            \filldraw (2.5,0) circle (2pt);
            \node[below=5pt] at (2.5,0) {$V_1 = \mathbb{C}x$};
            \node[above=5pt] at (2.5,0) {$1$};
            \draw[->, bend left=25] (-2.35,0.08) to node[above] {\small$E$} (2.35,0.08);
            \draw[->, bend left=25] (2.35,-0.08) to node[below] {\small$F$} (-2.35,-0.08);
        \end{tikzpicture}
        \end{center}

        \item $\mathfrak{sl}(2,\mathbb{C})$ acts on itself ($\alpha=2$) by
        \[
            \rho : \mathfrak{sl}(2,\mathbb{C}) \to \mathrm{End}(\mathfrak{sl}(2,\mathbb{C})),
            \qquad X \mapsto [X,\,\cdot\,].
        \]
        The weight spaces are $V_{-2}=\mathbb{C}F$, $V_0=\mathbb{C}H$, $V_2=\mathbb{C}E$:
        \begin{center}
        \begin{tikzpicture}[>=Stealth, baseline=(current bounding box.center)]
            \filldraw (-3,0) circle (2pt);
            \node[above=5pt] at (-3,0) {$F$};
            \node[below=5pt] at (-3,0) {$V_{-2}$};
            \filldraw (0,0) circle (2pt);
            \node[above=5pt] at (0,0) {$H$};
            \node[below=5pt] at (0,0) {$V_0$};
            \filldraw (3,0) circle (2pt);
            \node[above=5pt] at (3,0) {$E$};
            \node[below=5pt] at (3,0) {$V_2$};
            \draw[->, bend left=25] (-2.85,0.08) to node[above] {\small$E$} (-0.15,0.08);
            \draw[->, bend left=25] (-0.15,-0.08) to node[below] {\small$F$} (-2.85,-0.08);
            \draw[->, bend left=25] (0.15,0.08) to node[above] {\small$E$} (2.85,0.08);
            \draw[->, bend left=25] (2.85,-0.08) to node[below] {\small$F$} (0.15,-0.08);
        \end{tikzpicture}
        \end{center}
        This is called the \textbf{adjoint representation}\index{adjoint representation!of a Lie algebra}.
    \end{itemize}
}

\fact{
    All finite-dimensional $\mathfrak{sl}(2,\mathbb{C})$ representations are fully reducible.
}

\defn[tensor product representation!Lie algebra]{Tensor product representation}{
    Let $V$ and $W$ be representations of a Lie algebra $\mathfrak{g}$.
    The \textbf{tensor product representation}
    on $V\otimes W$ is defined by
    \[
        \rho_{V\otimes W}(X)
        :=
        \rho_V(X)\otimes\mathrm{id}_W + \mathrm{id}_V\otimes\rho_W(X)
        \qquad \forall\, X\in\mathfrak{g}.
    \]
    (Verification that this is a representation: see exercise sheets.)
}

\ex{
    \textbf{Isospin (Heisenberg, 1932).}\quad When only the strong interaction is
    considered, there is an approximate symmetry exchanging proton and neutron. This makes
    $\mathbb{C}[p,n]$ a copy of the standard representation (highest weight $\alpha=1$),
    with the proton at weight $+1$ and the neutron at weight $-1$. Extending to the three
    light quarks,
    \[
        \mathbb{C}[u,d,s] \cong V_1 \oplus V_0,
    \]
    with $u,d$ at weights $\pm 1$ and $s$ at weight $0$, where $V_n$ denotes the irreducible
    representation with highest weight $n$. Mesons transform in the following multiplets:
    \begin{itemize}
        \item $\pi^+, \pi^0, \pi^-$ form an isospin triplet in $V_2$;
        \item $(K^+, K^0)$ and $(K^-,\bar{K}^0)$ form isospin doublets in $V_1$.
    \end{itemize}

    \medskip\noindent\textit{Tensor product $V_1\otimes V_1$.}\quad
    Weights add under tensor products:
    $\rho(H)(v\otimes w) = (\lambda_v+\lambda_w)\,v\otimes w$.
    For $V_1$ with basis $\{x,y\}$ at weights $+1,-1$:
    \begin{center}
    \begin{tikzpicture}[>=Stealth, baseline=(current bounding box.center)]
        \filldraw (-3,0) circle (2pt);
        \node[above=5pt] at (-3,0) {$y\otimes y$};
        \node[below=5pt] at (-3,0) {$-2$};
        \filldraw (-0.15,0) circle (2pt);
        \filldraw ( 0.15,0) circle (2pt);
        \node[above=7pt] at (0,0) {$y\otimes x,\quad x\otimes y$};
        \node[below=5pt] at (0,0) {$0$};
        \filldraw (3,0) circle (2pt);
        \node[above=5pt] at (3,0) {$x\otimes x$};
        \node[below=5pt] at (3,0) {$2$};
    \end{tikzpicture}
    \end{center}
    The weight-$0$ space is $2$-dimensional, so $V_1\otimes V_1 \cong V_2\oplus V_0$.

    \medskip\noindent\textit{Tensor product $V_2\otimes V_1$.}\quad
    Let $\{|\mathord{\uparrow}\rangle, |0\rangle, |\mathord{\downarrow}\rangle\}$ be a basis
    of $V_2$ at weights $+2,0,-2$. The six tensor products and their weights are:
    \[
        \begin{array}{c|cc}
             & x\;(+1) & y\;(-1) \\
            \hline
            |\mathord{\uparrow}\rangle\;(+2)
                & |\mathord{\uparrow}\rangle\otimes x\;(+3)
                & |\mathord{\uparrow}\rangle\otimes y\;(+1) \\[4pt]
            |0\rangle\;(0)
                & |0\rangle\otimes x\;(+1)
                & |0\rangle\otimes y\;(-1) \\[4pt]
            |\mathord{\downarrow}\rangle\;(-2)
                & |\mathord{\downarrow}\rangle\otimes x\;(-1)
                & |\mathord{\downarrow}\rangle\otimes y\;(-3)
        \end{array}
    \]
    The weight spaces have dimensions $1,2,2,1$ at weights $+3,+1,-1,-3$, consistent
    with $V_2\otimes V_1 \cong V_3\oplus V_1$.
}

\thmp{Classification of irreducible representations of $\mathfrak{sl}(2,\mathbb{C})$}{
    The irreducible complex representations of $\mathfrak{sl}(2,\mathbb{C})$ are, up to
    isomorphism, exactly $\mathrm{Sym}^n(\mathbb{C}^2)$ for $n\in\mathbb{N}_0$, where
    $\mathbb{C}^2$ carries the standard representation $X\mapsto X$. In particular,
    $\mathrm{Sym}^n(\mathbb{C}^2)$ is $(n+1)$-dimensional with highest weight $n$.
}{
    Let $x,y$ be the standard basis of $\mathbb{C}^2$. The space $\mathrm{Sym}^n(\mathbb{C}^2)$
    is $(n+1)$-dimensional with basis $\{x^{n-k}y^k\}_{k=0}^n$. Using the Leibniz rule
    together with $\rho(H)x = x$ and $\rho(H)y = -y$:
    \[
        \rho(H)(x^{n-k}y^k)
        =
        (n-k)(+1)\,x^{n-k}y^k + k(-1)\,x^{n-k}y^k
        =
        (n-2k)\,x^{n-k}y^k.
    \]
    The highest weight is $n$, achieved by $x^n$. Since $\rho(E)x = 0$ in the standard
    representation, $\rho(E)(x^n) = 0$, so $x^n$ is a highest weight vector. Using
    $\rho(F)x = y$ and $\rho(F)y = 0$, repeated application gives
    \[
        \rho(F)^k(x^n) = \frac{n!}{(n-k)!}\,x^{n-k}y^k \neq 0
        \qquad \text{for } k = 0,1,\ldots,n.
    \]
    By the previous proposition, $\mathrm{Sym}^n(\mathbb{C}^2)$ is generated by $x^n$ and
    is therefore irreducible.

    Conversely, every finite-dimensional irreducible representation has highest weight
    $\alpha\in\mathbb{N}_0$ and dimension $\alpha+1$ (by the corollary). Since
    $\mathrm{Sym}^\alpha(\mathbb{C}^2)$ is irreducible of dimension $\alpha+1$ with highest
    weight $\alpha$, it is the unique such representation up to isomorphism.
}

\subsection{Unitary representations and the physics convention}

\rmkb{
    In physics, one is interested in \textbf{unitary representations}\index{unitary representation},
    i.e.\ group homomorphisms $G \to U(n)$. At the Lie algebra level, this means maps into
    \[
        \mathfrak{u}(n)
        = \bigl\{A \in M_n(\mathbb{C}) \bigm| A^\dagger = -A\bigr\},
    \]
    so every generator acts as a skew-adjoint operator. This condition is not preserved by
    complexification $(\mathord{\otimes}_\mathbb{R}\mathbb{C})$: the complexification introduces
    new elements for which no natural adjointness condition is imposed.

    Both $\mathfrak{su}(2)$ and $\mathfrak{sl}(2,\mathbb{R})$ are real forms of
    $\mathfrak{sl}(2,\mathbb{C})$, but their unitary representation theories are opposite in
    character:
    \begin{itemize}
        \item $\mathfrak{su}(2)$: all finite-dimensional representations can be made unitary
              (the associated Lie group $SU(2)$ is compact).
        \item $\mathfrak{sl}(2,\mathbb{R})$: there are no non-trivial finite-dimensional
              unitary representations.
    \end{itemize}

    The standard basis of $\mathfrak{su}(2)$, expressed in terms of $H, E, F$, is
    \[
        Z = iH = \begin{pmatrix}i & 0 \\ 0 & -i\end{pmatrix},
        \qquad
        Y = E-F = \begin{pmatrix}0 & 1 \\ -1 & 0\end{pmatrix},
        \qquad
        X = i(E+F) = \begin{pmatrix}0 & i \\ i & 0\end{pmatrix}.
    \]
    One checks that each satisfies $A^\dagger = -A$. In any unitary $\mathfrak{su}(2)$-%
    representation, the skew-adjointness of $Z = iH$ forces $\rho(H)$ to be self-adjoint,
    while the skew-adjointness of $Y$ and $X$ together give
    \[
        \rho(E)^\dagger = \rho(F).
    \]
}

\ex{
    The standard representation on $\mathbb{C}^2$ is unitary: $SU(2) \subset GL(2,\mathbb{C})$
    acts by isometries of the standard Hermitian inner product. Since tensor products of
    unitary representations are again unitary, $\mathrm{Sym}^n(\mathbb{C}^2)$ is unitary
    for all $n \in \mathbb{N}_0$.
}

\rmkb{
    \textbf{Convention.}\quad For a representation $\rho:\mathfrak{g}\to\mathrm{End}(V)$ we
    write $Xv := \rho(X)(v)$ for $X \in \mathfrak{g}$ and $v \in V$, suppressing $\rho$ from
    the notation.

    \medskip
    The basis $\{F^k v\}_{k=0}^n$ of $V_n$ is orthogonal (the vectors lie in distinct
    $\rho(H)$-eigenspaces, and $\rho(H)$ is self-adjoint), but not orthonormal. Using
    $\rho(E)^\dagger = \rho(F)$, which gives $\langle u \mid Fw\rangle = \langle Eu \mid w\rangle$,
    together with $EF^k v = k(n-k+1)F^{k-1}v$, one computes
    \begin{align*}
        \|F^k v\|^2
        &= \langle F^k v \mid F\cdot F^{k-1}v\rangle
         = \langle EF^k v \mid F^{k-1}v\rangle
         = k(n-k+1)\,\|F^{k-1}v\|^2.
    \end{align*}
    Starting from $\|v\|^2 = 1$, this recursion gives
    $\|F^k v\|^2 = k!\,\dfrac{n!}{(n-k)!}$
    and defines a positive-definite inner product with respect to which the representation is
    unitary.

    \medskip
    \textit{Physics basis.}\quad In quantum mechanics, states are labelled by eigenvalues of
    simultaneously measurable observables. Setting $j := n/2$ (total angular momentum quantum
    number), the orthonormal basis is defined inductively by
    \[
        \bigl|j,\, j\bigr\rangle := v,
        \qquad
        \bigl|j,\, m-1\bigr\rangle
        := \frac{F\,\bigl|j,m\bigr\rangle}{\sqrt{(j+m)(j-m+1)}},
        \qquad m = -j+1,\ldots,j.
    \]
    The second label $m$ (magnetic quantum number, or z-projection of angular momentum) runs
    over $m = -j,\,-j+1,\,\ldots,\,j$. In this basis the generators act by the standard
    angular momentum formulas:
    \begin{align*}
        H\,\bigl|j,m\bigr\rangle &= 2m\;\bigl|j,m\bigr\rangle,\\[4pt]
        E\,\bigl|j,m\bigr\rangle &= \sqrt{(j-m)(j+m+1)}\;\bigl|j,\,m+1\bigr\rangle.
    \end{align*}
}

\section{Representation theory of $\mathfrak{su}(3)$ and $\mathfrak{sl}(3,\mathbb{C})$}

\rmkb{
    \textbf{Physics background: the eightfold way.}\quad The three lightest quarks ---
    up ($u$), down ($d$), and strange ($s$) --- have masses
    \[
        m_u \approx 2.2\;\mathrm{MeV},\qquad
        m_d \approx 4.7\;\mathrm{MeV},\qquad
        m_s \approx 93\;\mathrm{MeV}.
    \]
    All three are much lighter than the typical hadronic scale (${\sim}1\;\mathrm{GeV}$),
    giving rise to an approximate $SU(3)$-flavor symmetry acting on $\mathbb{C}[u,d,s]$.
    Historically, physicists observed that mesons and baryons organize into multiplets matching
    representations of $\mathfrak{sl}(3)$ --- a phenomenon Gell-Mann called the
    \textbf{eightfold way}, named after the 8-dimensional adjoint representation in which the
    light mesons sit. Some of the predicted particles were initially absent from experiment;
    quarks were introduced as a mathematical bookkeeping device to explain the multiplet
    structure, and the missing particles were subsequently discovered. This led to Gell-Mann's
    Nobel Prize in 1969.
}

\defn[Cartan subalgebra]{Cartan subalgebra of $\mathfrak{sl}(3,\mathbb{C})$}{
    The \textbf{Cartan subalgebra} of $\mathfrak{sl}(3,\mathbb{C})$ is
    \[
        \mathfrak{h} :=
        \left\{
            \begin{pmatrix}a & 0 & 0 \\ 0 & b & 0 \\ 0 & 0 & c\end{pmatrix}
            \;\middle|\;
            a + b + c = 0
        \right\}
        \subseteq \mathfrak{sl}(3,\mathbb{C}).
    \]
    Since diagonal matrices commute, $\mathfrak{h}$ is abelian: $[X,Y]=0$ for all
    $X,Y\in\mathfrak{h}$. A basis is given by
    \[
        H_{12} = \begin{pmatrix}1&0&0\\0&-1&0\\0&0&0\end{pmatrix},
        \qquad
        H_{23} = \begin{pmatrix}0&0&0\\0&1&0\\0&0&-1\end{pmatrix}.
    \]
    The role of $\mathfrak{h}$ is analogous to that of $\mathbb{C} H$ in $\mathfrak{sl}(2)$:
    it is the subalgebra whose simultaneous eigenvectors we use to organise representations.
}

\fact{
    In every finite-dimensional $\mathfrak{sl}(3,\mathbb{C})$-representation, all elements
    of $\mathfrak{h}$ act by simultaneously diagonalizable operators.
}

\defn[-]{Weight}{
    Let $V$ be a finite-dimensional $\mathfrak{sl}(3,\mathbb{C})$-representation. A vector
    $v \in V$ is a \textbf{weight vector}\index{weight vector!$\mathfrak{sl}(3)$} with
    \textbf{weight} $\alpha \in \mathfrak{h}^*$\index{weight!$\mathfrak{sl}(3)$} if
    \[
        \rho(H)\,v = \alpha(H)\,v \qquad \forall\, H \in \mathfrak{h}.
    \]
    The corresponding \textbf{weight space}\index{weight space} is
    $V_\alpha = \{v \in V \mid \rho(H)v = \alpha(H)v\ \forall\, H\in\mathfrak{h}\}$.
    By the fact above, $V = \bigoplus_\alpha V_\alpha$ (finite direct sum over $\mathfrak{h}^*$).
}

\rmkb{
    Define linear functionals $L_1, L_2, L_3 \in \mathfrak{h}^*$ by
    \[
        L_i\!\begin{pmatrix}a_1&0&0\\0&a_2&0\\0&0&a_3\end{pmatrix} := a_i.
    \]
    Since every $H \in \mathfrak{h}$ satisfies $a_1+a_2+a_3=0$, we have
    $L_1+L_2+L_3=0$, so $\{L_1, L_2\}$ is a basis of $\mathfrak{h}^*$ with
    $L_3 = -L_1-L_2$. The three functionals are symmetric and point in directions
    $120^\circ$ apart in $\mathfrak{h}^*$:
    \begin{center}
    \begin{tikzpicture}[>=Stealth, scale=1.1, baseline=(current bounding box.center)]
        \draw[->] (0,0) -- ({cos(-30)},{sin(-30)}) node[right] {$L_1$};
        \draw[->] (0,0) -- ({cos(90)}, {sin(90)})  node[above] {$L_2$};
        \draw[->] (0,0) -- ({cos(210)},{sin(210)}) node[left]  {$L_3$};
        \filldraw (0,0) circle (1.5pt);
    \end{tikzpicture}
    \end{center}
}

\ex{
    The defining 3-dimensional representation $\mathfrak{sl}(3,\mathbb{C}) \to \mathrm{End}(\mathbb{C}^3)$,
    $X \mapsto X$, decomposes into weight spaces as
    \[
        \mathbb{C}^3 = \mathbb{C}_{L_1} \oplus \mathbb{C}_{L_2} \oplus \mathbb{C}_{L_3},
    \]
    where $e_i$ has weight $L_i$: $H e_i = a_i e_i = L_i(H)\,e_i$ for
    $H = \mathrm{diag}(a_1,a_2,a_3)$.
}

\rmkb{
    In $\mathfrak{sl}(2)$ we found eigenvectors of $\mathrm{ad}(H)$ (i.e.\ $[H,E]=2E$,
    $[H,F]=-2F$) to build a basis that moves between weight spaces. We do the same for
    $\mathfrak{sl}(3)$: decompose $\mathfrak{sl}(3,\mathbb{C})$ into weight spaces under
    the adjoint representation $X \mapsto [X,-]$.

    Since $[\mathfrak{h}, \mathfrak{h}] = 0$, the Cartan subalgebra sits at weight $0$.
    The remaining basis elements are the \textbf{matrix units} $E_{ij}$ (with $1$ in
    position $(i,j)$ and $0$ elsewhere, $i \neq j$). For $H = \mathrm{diag}(a_1,a_2,a_3)$:
    \begin{align*}
        [H, E_{12}]
        &= \begin{pmatrix}a_1&0&0\\0&a_2&0\\0&0&a_3\end{pmatrix}
           \begin{pmatrix}0&1&0\\0&0&0\\0&0&0\end{pmatrix}
           -
           \begin{pmatrix}0&1&0\\0&0&0\\0&0&0\end{pmatrix}
           \begin{pmatrix}a_1&0&0\\0&a_2&0\\0&0&a_3\end{pmatrix}\\[4pt]
        &= \begin{pmatrix}0&a_1&0\\0&0&0\\0&0&0\end{pmatrix}
           -
           \begin{pmatrix}0&a_2&0\\0&0&0\\0&0&0\end{pmatrix}
           = (a_1 - a_2)\,E_{12}
           = (L_1-L_2)(H)\,E_{12}.
    \end{align*}
    The same computation gives $[H, E_{ij}] = (L_i-L_j)(H)\,E_{ij}$ for all $H \in \mathfrak{h}$,
    so $E_{ij}$ is a weight vector with weight $L_i-L_j$. This gives the
    \textbf{root space decomposition}\index{root space decomposition}
    \[
        \mathfrak{sl}(3,\mathbb{C}) = \mathfrak{h}
        \;\oplus\;
        \bigoplus_{i \neq j}\mathfrak{g}_{L_i-L_j},
        \qquad
        \mathfrak{g}_{L_i-L_j} = \mathbb{C}\,E_{ij}.
    \]
    The six weights $L_i-L_j$ ($i\neq j$), called \textbf{roots}\index{root}, form a regular
    hexagon in $\mathfrak{h}^*$. The total dimension $\dim\mathfrak{sl}(3,\mathbb{C}) = 2+6 = 8$
    matches the eightfold way.
    \begin{center}
    \begin{tikzpicture}[>=Stealth, scale=1.8, baseline=(current bounding box.center)]
        \coordinate (UR) at (60:1);
        \coordinate (UL) at (120:1);
        \coordinate (L)  at (180:1);
        \coordinate (LL) at (240:1);
        \coordinate (LR) at (300:1);
        \coordinate (R)  at (0:1);
        \filldraw[fill=gray!15, draw=black] (UR)--(UL)--(L)--(LL)--(LR)--(R)--cycle;
        \foreach \v in {UR,UL,L,LL,LR,R} { \filldraw (\v) circle (1.5pt); }
        \filldraw (0,0) circle (1.5pt);
        \node[above right=1pt] at (UR) {$L_2{-}L_3$};
        \node[above left=1pt]  at (UL) {$L_2{-}L_1$};
        \node[left=1pt]        at (L)  {$L_3{-}L_1$};
        \node[below left=1pt]  at (LL) {$L_3{-}L_2$};
        \node[below right=1pt] at (LR) {$L_1{-}L_2$};
        \node[right=1pt]       at (R)  {$L_1{-}L_3$};
        \draw[->] (0,0) -- (330:0.55) node[right=1pt] {\small$L_1$};
        \draw[->] (0,0) -- (90:0.55)  node[above=1pt] {\small$L_2$};
        \draw[->] (0,0) -- (210:0.55) node[left=1pt]  {\small$L_3$};
    \end{tikzpicture}
    \end{center}
}

\propp{
    Let $X, Y \in \mathfrak{sl}(3,\mathbb{C})$ be weight vectors of weights $\alpha$ and
    $\beta$. Then $[X,Y]$ has weight $\alpha+\beta$, i.e.\
    \[
        [X,Y] \;\in\; \mathfrak{g}_{\alpha+\beta},
    \]
    where $\mathfrak{g}_\gamma := 0$ whenever $\gamma$ is not a root.
}{
    For any $H \in \mathfrak{h}$, the Jacobi identity gives
    \begin{align*}
        [H, [X,Y]]
            &= -[Y, [H,X]] - [X, [Y,H]] \\
            &= -[Y, \alpha(H)X] + [X, \beta(H)Y] \\
            &= \bigl(\alpha(H) + \beta(H)\bigr)\,[X,Y].
    \end{align*}
    Since this holds for all $H \in \mathfrak{h}$, the element $[X,Y]$ is a weight
    vector of weight $\alpha+\beta$.
}

\ex{
    In the diagram, taking the commutator with $E_{ij}$ moves every weight in the
    $E_{ij}$-direction. The diagram below shows the action of $[E_{31}, -]$, which
    shifts by $L_3 - L_1$, i.e.\ one step to the left. Moving outside the hexagon
    gives zero.
    \begin{center}
    \begin{tikzpicture}[>=Stealth, scale=2.0, baseline=(current bounding box.center)]
        \coordinate (UR) at (60:1);
        \coordinate (UL) at (120:1);
        \coordinate (L)  at (180:1);
        \coordinate (LL) at (240:1);
        \coordinate (LR) at (300:1);
        \coordinate (R)  at (0:1);
        \coordinate (O)  at (0,0);
        % Hexagon outline
        \draw[gray!40] (UR)--(UL)--(L)--(LL)--(LR)--(R)--cycle;
        \foreach \v in {UR,UL,L,LL,LR,R} { \filldraw[black!50] (\v) circle (1.5pt); }
        \filldraw[black!30] (O) circle (1.5pt);
        % Root weight labels with the matching E_{ij} stacked directly underneath
        \node[above right=5pt] at (UR) {\shortstack{$L_2{-}L_3$\\[2pt]{\scriptsize\color{gray}$E_{23}$}}};
        \node[above left=5pt]  at (UL) {\shortstack{$L_2{-}L_1$\\[2pt]{\scriptsize\color{gray}$E_{21}$}}};
        \node[left=8pt]        at (L)  {\shortstack{$L_3{-}L_1$\\[2pt]{\scriptsize\color{gray}$E_{31}$}}};
        \node[below left=5pt]  at (LL) {\shortstack{$L_3{-}L_2$\\[2pt]{\scriptsize\color{gray}$E_{32}$}}};
        \node[below right=5pt] at (LR) {\shortstack{$L_1{-}L_2$\\[2pt]{\scriptsize\color{gray}$E_{12}$}}};
        \node[right=8pt]       at (R)  {\shortstack{$L_1{-}L_3$\\[2pt]{\scriptsize\color{gray}$E_{13}$}}};
        % Center label
        \node[above=3pt, font=\small, black!60] at (O) {$\mathfrak{h}$};
        % L-direction arrows (orientation guide)
        \draw[->, gray!45] (O) -- (330:0.50) node[right=2pt, font=\scriptsize, gray!65] {$L_1$};
        \draw[->, gray!45] (O) -- ( 90:0.50) node[above=1pt, font=\scriptsize, gray!65] {$L_2$};
        \draw[->, gray!45] (O) -- (210:0.50) node[left=2pt,  font=\scriptsize, gray!65] {$L_3$};
        % [E31, -] action arrows: shift by L3-L1 (leftward)
        \draw[->, blue, thick] ($(UR)!0.10!(UL)$) -- ($(UR)!0.90!(UL)$);  % E23 -> E21
        \draw[->, blue, thick] ($(LR)!0.10!(LL)$) -- ($(LR)!0.90!(LL)$);  % E12 -> E32
        \draw[->, blue!55, thick] ($(R)!0.12!(O)$)  -- ($(R)!0.88!(O)$);  % E13 -> h
        \draw[->, blue!55, thick] ($(O)!0.12!(L)$)  -- ($(O)!0.88!(L)$);  % h   -> E31
    \end{tikzpicture}
    \end{center}
}

\propp{
    Let $V$ be a finite-dimensional representation of $\mathfrak{sl}(3,\mathbb{C})$, let
    $v \in V$ be a weight vector of weight $\alpha$, and let $X \in \mathfrak{sl}(3,\mathbb{C})$
    be of weight $\beta$. Then $X(v) \in V$ is of weight $\alpha+\beta$.
}{
    For any $H \in \mathfrak{h}$,
    \[
        H(X(v)) = X(H(v)) + [H,X](v) = \bigl(\alpha+\beta\bigr)(H)\,X(v).
    \]
}

\rmkb{
    It follows that the weights appearing in an irreducible representation differ by
    integer combinations of the roots. The lattice generated by the roots,
    \[
        \Lambda_R \;\subseteq\; \mathfrak{h}^*,
    \]
    is called the \textbf{root lattice}\index{root lattice}.
}

There is no canonical notion of a highest weight, so we must make a choice.

\defn[highest weight map]{Highest weight map}{
    A \textbf{highest weight map} is a linear map $\ell \colon \mathfrak{h}^* \to \mathbb{C}$
    such that
    \[
        \operatorname{Re}\bigl(\ell(\alpha)\bigr) \neq 0
        \qquad \text{for all } 0 \neq \alpha \in \Lambda_R.
    \]
    A root $\alpha$ is called \textbf{positive}\index{root!positive} if $\ell(\alpha) > 0$,
    and \textbf{negative}\index{root!negative} otherwise.
}

To be concrete, for $\mathfrak{sl}(3)$ we pick
\[
    \ell(a_1 L_1 + a_2 L_2 + a_3 L_3) \;=\; a\,a_1 + b\,a_2 + c\,a_3,
    \qquad a + b + c = 0,\quad a > b > c.
\]
This particular choice does not matter.

\defn[-]{Highest weight vector ($\mathfrak{sl}(3)$)}{
    Let $V$ be a representation of $\mathfrak{sl}(3,\mathbb{C})$. A weight vector
    $v \in V_\alpha$ is called a \textbf{highest weight vector}\index{weight!highest weight vector}
    if $\operatorname{Re}\,\ell(\alpha)$ is maximal among all weights of $V$, or equivalently if
    the positive root vectors annihilate it:
    \[
        E_{12}\,v \;=\; E_{23}\,v \;=\; E_{13}\,v \;=\; 0.
    \]
}

\rmkb{
    The word \emph{weight} is used in several related senses, which are worth keeping apart:
    \begin{itemize}
        \item A \textbf{weight} is, abstractly, an element $\alpha \in \mathfrak{h}^*$ of the
              dual of the Cartan subalgebra --- a linear functional on $\mathfrak{h}$.
        \item The weight of a \textbf{vector} $v$ in a representation $\rho \colon
              \mathfrak{sl}(3,\mathbb{C}) \to \mathrm{End}(V)$ is the $\alpha \in \mathfrak{h}^*$
              with
              \[
                  \rho(H)\,v = \alpha(H)\,v \qquad \forall\, H \in \mathfrak{h};
              \]
              such a $v$ is a \emph{weight vector} of weight $\alpha$. (We usually omit $\rho$
              and write $Hv := \rho(H)v$.)
        \item The \textbf{weight space} $V_\alpha$ is the span of all weight vectors of weight
              $\alpha$. A weight of the \textbf{representation} $V$ is then an $\alpha$ with
              $V_\alpha \neq 0$, i.e.\ one that actually occurs in the decomposition
              $V = \bigoplus_\alpha V_\alpha$.
        \item In the special case of the \textbf{adjoint representation}, where
              $\mathfrak{sl}(3,\mathbb{C})$ acts on itself by $X \mapsto [X,-]$, the (non-zero)
              weights are exactly the \textbf{roots} $L_i - L_j$, and the corresponding weight
              vectors are the root vectors $E_{ij}$.
    \end{itemize}
}

\propp{
    Let $V$ be an irreducible finite-dimensional representation of $\mathfrak{sl}(3,\mathbb{C})$
    and $v \in V_\alpha$ a highest weight vector. Then $V$ is generated by the elements
    obtained from $v$ by successively applying $E_{21}$, $E_{31}$, and $E_{32}$.
}{
    Since $V$ is irreducible, it suffices to show that the subspace generated in this way
    is a non-zero subrepresentation; irreducibility then forces it to equal $V$.

    Let $W \leq V$ be the subspace generated by acting on $v$ with $E_{21}$, $E_{31}$,
    and $E_{32}$. We need to show that acting with the positive root vectors $E_{12}$,
    $E_{13}$, $E_{23}$ also preserves $W$. Since $E_{13} = [E_{12}, E_{23}]$, it is enough
    to show $E_{12}(W) \subseteq W$ and $E_{23}(W) \subseteq W$.

    Note that $W$ can equally be generated by acting with only $E_{21}$ and $E_{32}$, since
    $[E_{21}, E_{32}] = -E_{31}$.

    We proceed by induction. Define $W_n$ to be the subspace of $W$ generated by applying
    words of length at most $n$ in $E_{21}, E_{32}$ to $v$, with $W_{-1} := 0$ and
    $W_0 := \mathbb{C}v$. We show that acting with a positive root vector maps $W_n$ into
    $W_{n-1}$.

    Let $w_n \in W_n$. There are two cases,
    \[
        w_n = E_{21}\,w_{n-1} \qquad\text{or}\qquad w_n = E_{32}\,w_{n-1},
        \qquad w_{n-1} \in W_{n-1}.
    \]
    In the first case,
    \begin{align*}
        E_{12}\,E_{21}\,w_{n-1}
            &= E_{21}\,E_{12}\,w_{n-1} + \underbrace{[E_{12}, E_{21}]}_{\in\,\mathfrak{h}}\,w_{n-1} \\
            &= E_{21}\,\underbrace{E_{12}\,w_{n-1}}_{\in\,W_{n-2}} + \alpha\bigl([E_{12}, E_{21}]\bigr)\,w_{n-1}.
    \end{align*}
    By the induction hypothesis $E_{12}\,w_{n-1} \in W_{n-2}$, so the first term lies in
    $W_{n-1}$; the second term lies in $W_{n-1}$ as well. Hence $E_{12}\,w_n \in W_{n-1}$.

    The remaining cases follow in the same way from
    \[
        [E_{23}, E_{12}] = 0, \qquad [E_{12}, E_{32}] = 0, \qquad [E_{23}, E_{32}] \in \mathfrak{h}.
    \]
}

\rmkb{
    Let $V$ be a finite-dimensional representation of $\mathfrak{sl}(3,\mathbb{C})$ and
    $v \in V_\alpha$ a highest weight vector. Then acting with the negative roots on $v$
    generates an irreducible subrepresentation.

    The weights reachable from $v$ fill a cone opening in the directions of the negative
    roots. Along a single root string the picture reduces to a familiar
    $\mathfrak{sl}(2,\mathbb{C})$-representation.
    \begin{center}
    \begin{tikzpicture}[>=Stealth, scale=1.1, baseline=(current bounding box.center)]
        % negative root step vectors: E21 ~ 120 deg, E32 ~ 240 deg
        % lattice point (i,j) = i*(120 deg) + j*(240 deg), v = (0,0) at the right
        % Cone of possible weights
        \foreach \i in {0,1,2,3}{
            \pgfmathtruncatemacro{\jmax}{3-\i}
            \foreach \j in {0,...,\jmax}{
                \filldraw[black!45] ({-0.5*(\i+\j)},{0.866*(\i-\j)}) circle (1.3pt);
            }
        }
        % dashed enclosure around the E32-string (i=0 column)
        \draw[dashed, gray!70, rotate around={60:(-0.75,-1.3)}]
            (-0.75,-1.3) ellipse (1.95 and 0.42);
        % highest weight vector
        \filldraw (0,0) circle (1.8pt) node[right=3pt] {$v$};
        % negative-root arrows from v
        \draw[->, thick] (0,0) -- (-0.5, 0.866) node[midway, above left=-1pt, font=\small] {$E_{21}$};
        \draw[->, thick] (0,0) -- (-0.5,-0.866) node[midway, right=2pt, font=\small] {$E_{32}$};
        % labels
        \node[align=center, font=\small] at (-3.0,0.5) {Possible\\weights};
        \draw[->, gray] (-2.3,0.45) -- (-1.6,0.25);
        \node[font=\small] at (-0.2,-2.9) {Looks like an $\mathfrak{sl}(2)$-representation.};
        \draw[->, gray] (-0.6,-2.65) -- (-1.0,-2.0);
    \end{tikzpicture}
    \end{center}
}

\defn[Lie subalgebra]{Lie subalgebra}{
    Let $\mathfrak{g}$ be a Lie algebra. A \textbf{Lie subalgebra} is a subspace
    $\mathfrak{h} \leq \mathfrak{g}$ such that
    \[
        [X,Y] \in \mathfrak{h} \qquad \forall\, X,Y \in \mathfrak{h}.
    \]
}

\ex{
    For each pair $1 \leq i < j \leq 3$, the matrices $E_{ij}, E_{ji}, H_{ij}$ span a Lie
    subalgebra of $\mathfrak{sl}(3,\mathbb{C})$ isomorphic to $\mathfrak{sl}(2,\mathbb{C})$,
    with $E_{ij}, E_{ji}, H_{ij}$ taking the roles of $E, F, H$.

    Consequently the possible highest weights are labelled by two non-negative integers
    $(n_{12}, n_{23}) \in \mathbb{N}_0^2$, and the corresponding highest weight is
    \[
        (n_{12} + n_{23})\,L_1 + n_{23}\,L_2 \;=\; n_{12}\,L_1 - n_{23}\,L_3 .
    \]
}

In particular, every weight occurring in a finite-dimensional representation is an integer
combination $a L_1 + b L_2 + c L_3$ with $a, b, c \in \mathbb{Z}$.

\defn[weight lattice]{Weight lattice}{
    The \textbf{weight lattice} $\Lambda_W \subseteq \mathfrak{h}^*$ is the lattice generated
    by $L_1, L_2, L_3$. It contains the root lattice:
    \[
        \Lambda_R \;\subseteq\; \Lambda_W \;\subseteq\; \mathfrak{h}^*.
    \]
}

\rmkb{
    Restricting a representation to the $\mathfrak{sl}(2,\mathbb{C})$-subalgebra generated by
    $E_{ij}, E_{ji}, H_{ij}$ shows, exactly as in the $\mathfrak{sl}(2)$ theory, that its set of
    weights is symmetric under reflection in the hyperplane
    \[
        \{\lambda \in \mathfrak{h}^* \mid \lambda(H_{ij}) = 0\}.
    \]
    The group generated by these reflections is the \textbf{Weyl group}\index{Weyl group}.
    Every weight diagram is invariant under it, so one completes a weight diagram by reflecting
    its weights in all the mirror hyperplanes.
    \begin{center}
    \begin{tikzpicture}[>=Stealth, scale=1.5, baseline=(current bounding box.center)]
        % three mirror hyperplanes lambda(H_ij)=0, at 30/90/150 degrees (mod 180)
        \foreach \m in {30,90,150}{
            \draw[dashed, gray!55] (\m+180:1.95) -- (\m:1.95);
        }
        % L_i direction arrows (lie along the mirror lines)
        \draw[->, gray!50] (0,0) -- (330:0.6) node[font=\scriptsize, gray!70, below right=-2pt] {$L_1$};
        \draw[->, gray!50] (0,0) -- ( 90:0.6) node[font=\scriptsize, gray!70, right=1pt] {$L_2$};
        \draw[->, gray!50] (0,0) -- (210:0.6) node[font=\scriptsize, gray!70, below left=-2pt] {$L_3$};
        % Weyl orbit of a dominant weight (regular hexagon)
        \foreach \a in {0,60,120,180,240,300}{
            \filldraw[black!45] (\a:1.5) circle (2pt);
        }
        % highest (dominant) weight, in the fundamental chamber between mirrors 30 and 90
        \filldraw[red!65!black] (60:1.5) circle (2.5pt) node[above right=0pt, font=\small] {$\lambda$};
        % reflection of lambda across the mirror at 30 deg
        \draw[->, red!50!black, bend right=20] (54:1.5) to node[right=1pt, font=\scriptsize] {} (6:1.5);
        \filldraw (0,0) circle (1pt);
        % mirror label
        \node[font=\small] at (30:2.35) {$\lambda(H_{12})=0$};
    \end{tikzpicture}
    \end{center}
}

Recall the approximate $SU(3)$-flavor symmetry rotating the light quarks $u, d, s$. One might
ask why the relevant symmetry is the continuous group $SU(3)$ rather than a discrete group such
as $\mathbb{Z}_3$; the answer is that the observed hadron multiplets match precisely the
irreducible representations of $\mathfrak{sl}(3,\mathbb{C})$, as the following examples
illustrate.

\ex{
    \textbf{Baryon octet (adjoint representation, $\mathbf{8}$).}\quad
    The spin-$\tfrac12$ baryons fill the weight diagram of the adjoint representation, a
    hexagon with a doubly-occupied centre:
    \begin{center}
    \begin{tikzpicture}[>=Stealth, scale=1.5, baseline=(current bounding box.center)]
        \draw[gray!35] (-0.5,1)--(0.5,1)--(1,0)--(0.5,-1)--(-0.5,-1)--(-1,0)--cycle;
        % blue: isospin (horizontal) steps
        \draw[->, blue!70, shorten >=3pt, shorten <=3pt] (-0.5,1) -- (0.5,1);
        \draw[->, blue!70, shorten >=3pt, shorten <=3pt] (-1,0) -- (-0.15,0);
        \draw[->, blue!70, shorten >=3pt, shorten <=3pt] (0.15,0) -- (1,0);
        \draw[->, blue!70, shorten >=3pt, shorten <=3pt] (-0.5,-1) -- (0.5,-1);
        % red: U-spin (diagonal) steps
        \draw[->, red!70, shorten >=3pt, shorten <=3pt] (-1,0) -- (-0.5,1);
        \draw[->, red!70, shorten >=3pt, shorten <=3pt] (-0.5,-1) -- (0,0);
        \draw[->, red!70, shorten >=3pt, shorten <=3pt] (0,0) -- (0.5,1);
        \draw[->, red!70, shorten >=3pt, shorten <=3pt] (0.5,-1) -- (1,0);
        % particles
        \filldraw (-0.5,1) circle (2pt) node[above=2pt] {$n$};
        \filldraw (0.5,1)  circle (2pt) node[above=2pt] {$p$};
        \filldraw (-1,0)   circle (2pt) node[left=2pt]  {$\Sigma^-$};
        \filldraw (1,0)    circle (2pt) node[right=2pt] {$\Sigma^+$};
        \filldraw (-0.09,0) circle (2pt);
        \filldraw (0.09,0)  circle (2pt);
        \node[below right=1pt] at (0.09,0) {$\Sigma^0,\Lambda$};
        \filldraw (-0.5,-1) circle (2pt) node[below=2pt] {$\Xi^-$};
        \filldraw (0.5,-1)  circle (2pt) node[below=2pt] {$\Xi^0$};
    \end{tikzpicture}
    \end{center}

    \medskip
    \textbf{Baryon decuplet ($\mathbf{10}$).}\quad
    The spin-$\tfrac32$ baryons fill a triangular weight diagram. The bottom vertex
    $\Omega^-$ was predicted before being discovered experimentally in $1964$:
    \begin{center}
    \begin{tikzpicture}[>=Stealth, scale=1.2, baseline=(current bounding box.center)]
        \draw[gray!35] (-1.5,1)--(1.5,1)--(0,-2)--cycle;
        \filldraw (-1.5,1) circle (2pt) node[above=2pt] {$\Delta^-$};
        \filldraw (-0.5,1) circle (2pt) node[above=2pt] {$\Delta^0$};
        \filldraw (0.5,1)  circle (2pt) node[above=2pt] {$\Delta^+$};
        \filldraw (1.5,1)  circle (2pt) node[above=2pt] {$\Delta^{++}$};
        \filldraw (-1,0)   circle (2pt) node[left=2pt]  {$\Sigma^{*-}$};
        \filldraw (0,0)    circle (2pt) node[below right=0pt] {$\Sigma^{*0}$};
        \filldraw (1,0)    circle (2pt) node[right=2pt] {$\Sigma^{*+}$};
        \filldraw (-0.5,-1) circle (2pt) node[left=2pt]  {$\Xi^{*-}$};
        \filldraw (0.5,-1)  circle (2pt) node[right=2pt] {$\Xi^{*0}$};
        \filldraw (0,-2)    circle (2pt) node[below=2pt] {$\Omega^-$};
    \end{tikzpicture}
    \end{center}

    \medskip
    \textbf{Quarks ($\mathbf{3}$) and antiquarks ($\bar{\mathbf{3}}$).}\quad
    The three quarks sit in the fundamental representation, the three antiquarks in its dual:
    \begin{center}
    \begin{tikzpicture}[>=Stealth, scale=1.3, baseline=(current bounding box.center)]
        \begin{scope}[xshift=-1.6cm]
            \draw[gray!35] (-0.6,0.4)--(0.6,0.4)--(0,-0.55)--cycle;
            \filldraw (-0.6,0.4) circle (2pt) node[above=1pt] {$d$};
            \filldraw (0.6,0.4)  circle (2pt) node[above=1pt] {$u$};
            \filldraw (0,-0.55)  circle (2pt) node[below=1pt] {$s$};
            \node at (0,-1.2) {$\mathbf{3}$};
        \end{scope}
        \begin{scope}[xshift=1.6cm]
            \draw[gray!35] (0,0.55)--(0.6,-0.4)--(-0.6,-0.4)--cycle;
            \filldraw (0,0.55)   circle (2pt) node[above=1pt] {$\bar{s}$};
            \filldraw (-0.6,-0.4) circle (2pt) node[below=1pt] {$\bar{u}$};
            \filldraw (0.6,-0.4)  circle (2pt) node[below=1pt] {$\bar{d}$};
            \node at (0,-1.2) {$\bar{\mathbf{3}}$};
        \end{scope}
    \end{tikzpicture}
    \end{center}

    \medskip
    \textbf{Tensor products.}\quad
    Recall that the tensor product representation acts by
    \[
        \rho_{V\otimes W}(X) = \mathrm{id}_V \otimes \rho_W(X) + \rho_V(X) \otimes \mathrm{id}_W .
    \]
    Combining a quark with an antiquark gives the meson multiplets,
    \[
        \mathbf{3} \otimes \bar{\mathbf{3}} = \mathbf{8} \oplus \mathbf{1},
    \]
    a hexagon (the octet) together with a singlet at the centre:
    \begin{center}
    \begin{tikzpicture}[>=Stealth, scale=1.1, baseline=(current bounding box.center)]
        \draw[gray!35] (-0.5,0.9)--(0.5,0.9)--(1,0)--(0.5,-0.9)--(-0.5,-0.9)--(-1,0)--cycle;
        \foreach \p in {(-0.5,0.9),(0.5,0.9),(1,0),(0.5,-0.9),(-0.5,-0.9),(-1,0)}{
            \filldraw \p circle (2pt);
        }
        \filldraw (0,0) circle (2pt);
        \draw (0,0) circle (5pt);
        \node[right=8pt] at (1,0) {$=\ \mathbf{8} \oplus \mathbf{1}$};
    \end{tikzpicture}
    \end{center}
    Combining two quarks gives
    \[
        \mathbf{3} \otimes \mathbf{3} = \mathbf{6} \oplus \bar{\mathbf{3}},
    \]
    where the symmetric sextet $\mathbf{6}$ consists of the diquark states
    \[
        uu,\quad ud,\quad dd,\quad us,\quad ds,\quad ss.
    \]
    These diquarks cannot exist as free particles in nature because of \emph{confinement}.

    \medskip
    Finally, $\mathbf{3} \otimes \mathbf{8} = \mathbf{15} \oplus \bar{\mathbf{6}} \oplus \mathbf{3}$
    (left as an exercise).
}

We found that irreducible representations are determined by their highest weight. It remains
to ask whether every allowed weight actually occurs as the highest weight of an irreducible
representation. This can also be shown using Schur functors, but we take a different route.

\thmp{Classification of irreducible representations of $\mathfrak{sl}(3,\mathbb{C})$}{
    For each highest weight $\alpha = n_1 L_1 - n_2 L_3$ with $n_1, n_2 \in \mathbb{N}_0$,
    there exists a unique (up to isomorphism) irreducible representation of
    $\mathfrak{sl}(3,\mathbb{C})$ with highest weight $\alpha$.
}{
    \textit{Existence.}\quad
    If $V$ and $W$ are representations with highest weights $\lambda_V$ and $\lambda_W$, then
    $V \otimes W$ has highest weight $\lambda_V + \lambda_W$ (the tensor product of the two
    highest weight vectors is again a highest weight vector). Since $\mathbf{3}$ has highest
    weight $L_1$ and $\bar{\mathbf{3}}$ has highest weight $-L_3$, the representation
    \[
        \mathbf{3}^{\otimes n_1} \otimes \bar{\mathbf{3}}^{\otimes n_2}
    \]
    has highest weight $n_1 L_1 - n_2 L_3$. Picking a highest weight vector $v$ in it,
    $v$ generates an irreducible subrepresentation with highest weight $n_1 L_1 - n_2 L_3$.
    This proves existence.

    \medskip
    \textit{Uniqueness.}\quad
    Let $V$ and $W$ be irreducible representations, both with highest weight $\alpha$, and let
    $v \in V$, $w \in W$ be highest weight vectors. The representation $V \oplus W$ is not
    irreducible, but $v + w \in V \oplus W$ is again a highest weight vector of weight $\alpha$.
    Let $\widetilde{W} \leq V \oplus W$ be the irreducible subrepresentation generated by
    $v + w$. The two projections
    \[
        \widetilde{W} \to V, \qquad \widetilde{W} \to W
    \]
    are non-zero homomorphisms of representations, so by Schur's lemma they are isomorphisms.
    Hence
    \[
        V \cong \widetilde{W} \cong W.
    \]
}

\rmkb{
\textit{(Added by transcriber, not part of the lecture.)}
Representations of a Lie group $H$ classify the particles of any gauge theory
with structure group $H$; the following makes this precise for the Standard Model \citep{FrancoisDifferentialgeometrygaugetheoryintroduction2021}.

\medskip
\textbf{Gauge groups.}\quad
A gauge theory with structure group $H$ is built on a principal $H$-bundle
$\pi\colon P\to M$ over spacetime $M$.
The \emph{gauge group}
\[
    \mathcal{H} \;=\; \mathrm{Aut}_v(P) \;:=\;
    \bigl\{\,\Phi\colon P\xrightarrow{\sim}P \;\bigm|\;
    \pi\circ\Phi=\pi,\;\Phi(ph)=\Phi(p)h\,\bigr\}
\]
is the group of fiber-preserving automorphisms; each element corresponds to an
equivariant map $\gamma\colon P\to H$ satisfying $\gamma(ph)=h^{-1}\gamma(p)h$,
and $\mathrm{Lie}\,\mathcal{H}\cong\Omega^0_{\mathrm{tens}}(P,\mathrm{Lie}\,H)$.
Incorporating spacetime diffeomorphisms gives the full symmetry group
\[
    \mathcal{G} \;=\; \mathrm{Diff}(M)\ltimes\mathcal{H},
\]
where $\mathrm{Diff}(M)$ acts on $\mathcal{H}$ by conjugation via bundle lifts.
Physical observables must be $\mathcal{H}$-invariant and covariant under $\mathrm{Diff}(M)$.
For the Standard Model, $H = SU(3)_c\times SU(2)_L\times U(1)_Y$.

\medskip
\textbf{Connections and $D^A$.}\quad
A \emph{connection} $A\in\Omega^1_{\mathrm{eq}}(P,\mathrm{Lie}\,H)$ satisfies
$\iota_{X^\#}A=X$ for all $X\in\mathrm{Lie}\,H$, splitting $TP=VP\oplus HP_A$
with horizontal projector $h_A\colon TP\to HP_A$.
The \emph{exterior covariant derivative} is the composition
\[
    D^A \;:=\; \mathrm{hor}_A\circ d
    \;:\;
    \Omega^k_{\mathrm{tens}}(P,V)
    \;\xrightarrow{\;d\;}\;
    \Omega^{k+1}(P,V)
    \;\xrightarrow{\;\mathrm{hor}_A\;}\;
    \Omega^{k+1}_{\mathrm{tens}}(P,V),
\]
where $(\mathrm{hor}_A\,\alpha)(X_1,\ldots,X_{k+1}):=\alpha(h_AX_1,\ldots,h_AX_{k+1})$.
On tensorial 0-forms, equivariance ($d\psi(X^\#)=-\rho_*(X)\psi$) gives
$D^A\psi = d\psi+\rho_*(A)\psi$.
The \emph{curvature}
\[
    F = dA+\tfrac{1}{2}[A,A]\;\in\;\Omega^2_{\mathrm{tens}}(P,\mathrm{Lie}\,H),
    \qquad D^AF=0\;\;\text{(Bianchi)},
\]
transforms in the adjoint representation of $H$ under $\mathcal{H}$:
\[
    A\;\longmapsto\;\gamma^{-1}A\gamma+\gamma^{-1}d\gamma,
    \qquad
    F\;\longmapsto\;\gamma^{-1}F\gamma.
\]

\medskip
\textbf{Gauge bosons.}\quad
Each generator of $\mathrm{Lie}\,H$ gives one gauge boson.
The decomposition $\mathrm{Lie}\,H=\mathfrak{su}(3)\oplus\mathfrak{su}(2)\oplus\mathfrak{u}(1)$
yields $8+3+1=12$: 8 gluons ($\mathfrak{su}(3)_c$), $W^\pm,W^3$ ($\mathfrak{su}(2)_L$),
and $B$ ($\mathfrak{u}(1)_Y$).
The Yang--Mills Lagrangian $L_{\mathrm{YM}}=\tfrac{1}{2}\mathrm{Tr}(F\wedge{*}F)$ is
$\mathcal{H}$-invariant; a mass term $m^2\mathrm{Tr}(A\wedge{*}A)$ is not, since $A$
transforms affinely---gauge invariance forces all gauge bosons massless.

\medskip
\textbf{Matter fields.}\quad
A representation $\rho\colon H\to GL(V)$ specifies a matter field: a tensorial 0-form
$\psi\in\Omega^0_{\mathrm{tens}}(P,V)\cong\Gamma(P\times_H V)$.
The \emph{gauge charge} of $\psi$ is the representation $\rho$; minimal coupling to
$A$ is automatic via $D^A$, and $\langle\psi,\slashed{D}^A\psi\rangle$ is
$\mathcal{H}$-invariant by equivariance.
For the quarks of each generation:
\[
    Q_L\sim(\mathbf{3},\mathbf{2},+\tfrac{1}{6}),\quad
    u_R\sim(\mathbf{3},\mathbf{1},+\tfrac{2}{3}),\quad
    d_R\sim(\mathbf{3},\mathbf{1},-\tfrac{1}{3}).
\]
The representation content is experimental input, constrained by
\emph{anomaly cancellation}: certain traces over the fermion representations must vanish.

\medskip
\textbf{Mass generation.}\quad
The Higgs field $\phi\in\Omega^0_{\mathrm{tens}}(P,\mathbb{C}^2)$ in representation
$(\mathbf{1},\mathbf{2},+\tfrac{1}{2})$ of $H$ couples to $A$ via
\[
    L_{\mathrm{Higgs}}=\langle D^A\phi,\,{*}D^A\phi\rangle-V(\phi),
    \qquad
    V(\phi)=\mu^2|\phi|^2+\lambda|\phi|^4.
\]
The full Lagrangian $L_{\mathrm{YM}}+L_{\mathrm{Higgs}}
+\sum_\psi\langle\psi,\slashed{D}^A\psi\rangle+L_{\mathrm{Yukawa}}$
is $\mathcal{H}$-invariant and entirely determined by $H$ and the representations
$\{\rho_i\}$; all couplings follow from $D^A$, and $L_{\mathrm{Yukawa}}$ generates
fermion masses.
Gauge boson mass generation admits two equivalent descriptions.

\smallskip
\textit{Spontaneous symmetry breaking.}\enspace
For $\mu^2<0$, a nonzero vacuum $\phi_0$ has stabilizer $U(1)_{\mathrm{EM}}\subset H$;
the residual gauge group is $SU(3)_c\times U(1)_{\mathrm{EM}}$.
The $W^\pm$ and $Z$ acquire masses from $|D^A\phi_0|^2$, with the Goldstone modes
absorbed as longitudinal polarizations; the photon remains massless.

\smallskip
\textit{Dressing field method} \citep{Attarddressingfieldmethodgaugesymmetryreductionreviewexamples2017}.\enspace
A \textbf{dressing field} for a subgroup $K\subset H$ is a $K$-valued field
$u\colon P\to K$ transforming as $u\mapsto\kappa^{-1}u$ under
$\kappa\in\mathcal{K}\subset\mathcal{H}$ (not as a connection).
The composite fields
\[
    A^u \;:=\; u^{-1}Au+u^{-1}du,\qquad \psi^u \;:=\; u^{-1}\psi
\]
are $\mathcal{K}$-invariant by construction.
For $K=SU(2)_L$, the polar decomposition $\phi=u[\phi]\,|\phi|$ provides $u[\phi]$
as a dressing field; the $W^\pm,Z$ masses appear directly from $|D^A\phi|^2$
rewritten in terms of $A^u$, with $|\phi|$ the physical neutral Higgs.
No vacuum, no Goldstone bosons, and no broken symmetry are invoked.

\smallskip
\textit{Dressing is not gauge fixing.}\enspace
Gauge fixing imposes a constraint $\chi(A)=0$, selects one representative per
$\mathcal{H}$-orbit, and requires Faddeev-Popov ghosts.
Dressing imposes no constraint: $(A,\phi)\mapsto(A^u,\psi^u)$ is a change of
variables on the full field space, globally well-defined (no Gribov copies), with
all degrees of freedom dynamical.
Both pictures yield the same S-matrix observables; Elitzur's theorem---local gauge
symmetry cannot be spontaneously broken---is automatically respected.

\medskip
\textbf{Flavor vs.\ gauge symmetry.}\quad
The $SU(3)_c$ factor is an exact gauge symmetry: color is a conserved Noether charge
and gluons couple identically to all quark flavors.
The $SU(3)_{\mathrm{flavor}}$ symmetry studied in this chapter is not a gauge symmetry:
it is an accidental near-symmetry of the light-quark mass terms, acts on flavor indices,
has no associated connection or gauge bosons, and is explicitly broken by quark mass
differences.
}
